\section{结论}
在5~mV主阈值下，128个响应中96项正确、20项判据误读、8项临界、4项数值不确定，失败和缺失各为0与0项。域内留出的存储0/1岭回归RMSE分别为0.088294/0.094644~fC（电荷）和0.798415/0.838043~mV（裕量损失）。尺寸外推的存储0/1岭回归RMSE分别为0.068049/0.067155~fC（电荷）和0.519434/0.504929~mV（裕量损失）。
本次完成了限定自然语言映射、真实事件到器件源项传递、匹配参考定义、统一主矩阵、预测前冻结、独立留出/外推评分与受控扩展的一致证据链。原科学门槛、样本分母和全部失败保持不变。边界验证显示连续量稳定与分类稳定需要分别报告，不能以一次连续门槛通过清除全部数值不确定。

下一步器件研究需用具有出处的工艺和电学数据校准结构、实际写入、保持、感放及恢复，扩展能量、位置与统计规模。课程方面已形成可编辑报告、研究稿、PPT与只读演示材料；真实学生贡献、教师审核、试讲、答辩和考核由参与者完成并据实留证。
